\section{探索算法}\label{sec:exploration}

\subsection{形式化}\label{sec:exploration:formalization}

探索循环维护一棵树 $\Tree=(\Nodes, E)$,其中每个节点
$n\in\Nodes$ 代表一次研究步骤:构想修订、代码改动、基准运行、
分析等。节点携带

\begin{itemize}
  \item 离散状态
        $\statefn{n}\in\{\StateOpen,\allowbreak\StateEval,\allowbreak\StateDone,\allowbreak\StateFailed\}$,
  \item 取自 \texttt{NodeLabel} 枚举的类别\term{标签},
  \item 用于按轴评估指标的自由形式 dict
        \code{node.metrics},含独立标量
        \code{\_scientific\_score}(取值范围 $[0,1]$),
  \item 标志节点是否携带实测值(而非临时文本)的布尔
        \texttt{has\_real\_data},以及
  \item 选择提示与剪枝谓词使用的簿记字段 \texttt{depth}。
        ARI 不重试失败节点,因此不会暴露任何 \texttt{retry\_count}
        信号。
\end{itemize}

将按轴信号汇总为标量 scientific score 的工作委托给任何实现
\texttt{Evaluator} 协议的对象;ARI 不固定线性组合。该协议是 BFTS 与下游评分细则评估之间唯一的接缝,
保持搜索引擎对如何将轴归约为标量保持中立。

\paragraph{Run-loop 状态.}
外层循环每次迭代变换以下元组:
\[
  \state \;=\; (\Frontier,\;\Pending,\;\expandcount,\;\Hist),
\]
其中 $\Frontier\subseteq\Nodes$ 为可扩展的完成节点集合
(\StateDone{} / \StateFailed{}),$\Pending\subseteq\Nodes$ 为
处于 \StateOpen{} 等待执行的节点集合,$\expandcount:\Nodes\to\Naturals$
记录每个前沿节点被扩展过几次,$\Hist=(l_1,\ldots,l_t)$ 是已执
行标签的滑动窗(长度上限 $\histlen=20$,详见
\cref{sec:exploration:diversity})。初始状态为
$\state_0=(\emptyset,\{n_\mathrm{root}\},\mathbf{0},())$。

\paragraph{剪枝谓词.}
硬性探索预算由如下谓词表达:
\begin{multline}\label{eq:prune}
  \prunepred(n;\,|\Nodes|) \;=\;
    \indicator{|\Nodes| \geq B_N}
    \;\lor\;
    \indicator{\depth{n} \geq B_D}
    \\\;\lor\;
    \sterilepred(n),
\end{multline}
其中 $B_N = \texttt{config.max\_total\_nodes}$,$B_D =
\texttt{config.max\_depth}$,sterile 谓词
\begin{equation}\label{eq:sterile}
  \sterilepred(n) \;=\; \indicator{\,|\Delta_n| = 0\,}
\end{equation}
($\Delta_n$ 为 $n$ 新增、修改或删除的文件集合)
当节点在 agent loop 结束时相对父节点未写出任何新工件时为真
。run loop 计算该文件差分,并把结果作为布尔
\texttt{\_sterile} 标志记录到节点上;\texttt{should\_prune} 只
读取该标志与两个上限。调用方每次迭代传入活动的 $|\Nodes|$,因
此节点数有唯一真值源。步数与成本预算由调用侧的成本追
踪器单独强制,而不在 BFTS 引擎内。

\paragraph{Retire 谓词.}
在 $\prunepred$ 之上,run-loop 还应用更软的\emph{前沿 retire 谓
词} $\retirepred$,在父节点 $f$ 不再是最高产线索时将其从
$\Frontier$ 中移除:
\begin{equation}\label{eq:retire}
  \retirepred(f, c) \;=\;
    \underbrace{\indicator{\nodescore{c} > \nodescore{f}}}_{\text{规则 A}}
    \;\lor\;
    \underbrace{\indicator{\expandcount(f) \geq \maxexp}}_{\text{规则 B}},
\end{equation}
其中 $\maxexp = \texttt{config.max\_expansions\_per\_node}$(默
认 $4$)。规则 A 在子节点 $c$ 一旦超越 $f$ 时退役 $f$;规则 B
限制每个父节点的预算,从而展开搜索面。
$\retirepred$ \emph{不删除}节点,其历史仍保留在 $\Tree$ 中,仅
撤销其继续被扩展的资格。

外层循环单次迭代的阶段分解(内层扩展子循环、并行 ReAct 批次、
执行后 retire)见 \cref{fig:bfts-state};显式伪代码见
\cref{sec:exploration:run-loop} 的 \cref{alg:bfts-run-loop}。

\begin{figure}[t]
  \centering
  \resizebox{0.85\columnwidth}{!}{% F15: BFTS state-machine view of one outer run-loop iteration, set in the
% ARI house design language (_style.tex). Single top-down spine; the two
% iteration scopes are dashed ariloop enclosures — the inner expansion
% sub-loop (steps 1-3) and the outer parallel-run loop (steps 4-7) —
% joined through the white state pools they communicate over (frontier
% F at the top, pending P between them). Pale-green boxes are the
% predicate gates (Pi, sigma, rho); the vermilion edges (the single
% accent role of this figure) are the side effects fired when a gate
% predicate evaluates true; the dotted gray margin path is the reflux
% into the next outer iteration. Auxiliary state (e, H) stays in the
% caption, exactly as before; per-edge state-bump annotations were
% folded into the caption to keep annotation density low.
%
% style.tikzstyles is loaded by the preamble (inlined) and by the
% standalone wrapper; the refined ari* styles used below come from
% _style.tex in this directory (gate T2: all widths live in the style
% files, never inline).
%
% The build cwd differs per entry point, so resolve _style.tex from each
% known prefix: report/<lang>/ (the language PDFs \input figures as
% ../shared/figures/tikz/<fig>.tex), report/ (make figure / render_tikz),
% the repo root, then this directory itself (direct cwd or TEXINPUTS).
\InputIfFileExists{../shared/figures/tikz/_style.tex}{}{%
  \InputIfFileExists{shared/figures/tikz/_style.tex}{}{%
    \InputIfFileExists{report/shared/figures/tikz/_style.tex}{}{%
      \InputIfFileExists{_style.tex}{}{%
        \errmessage{F15: cannot locate figures/tikz/_style.tex}}}}}
\begin{tikzpicture}[aricanvas]
  %% Heading: the state tuple the run loop mutates.
  \node[ariphase] (hd) {\figlabel{F15.band.state}};
  \pic at ([xshift=-2.5mm]hd.west) {ariglyph flask};

  %% Spine. White boxes are shared state pools, peach boxes the
  %% experimentation-phase steps, pale-green boxes the predicate gates.
  \node[aribox, below=0.4cm of hd]                    (frontier) {\figlabel{F15.frontier}};
  \node[ariboxwide, fill=phaseexp, below=0.9cm of frontier] (pick) {\figlabel{F15.pick_best}};
  \node[aribox gate, below=of pick]                   (prune)    {\figlabel{F15.prune}};
  \node[aribox exp, below=of prune]                   (expand)   {\figlabel{F15.expand}};
  \node[aribox, below=0.8cm of expand]                (pending)  {\figlabel{F15.pending}};
  \node[ariboxwide, fill=phaseexp, below=0.9cm of pending]  (batch) {\figlabel{F15.select_run}};
  \node[ariboxwide, fill=phaseexp, below=of batch]    (react)    {\figlabel{F15.run_parallel}};
  \node[aribox gate, below=of react]                  (sterile)  {\figlabel{F15.sterile}};
  \node[aribox, below=of sterile]                     (append)   {\figlabel{F15.append_frontier}};
  \node[aribox gate, below=of append]                 (retire)   {\figlabel{F15.retire}};

  %% Gate side effects (fire when the predicate is true).
  \node[aribox, right=of prune]   (dropf) {\figlabel{F15.prune_drop}};
  \node[aribox, right=of sterile] (clamp) {\figlabel{F15.clamp}};
  \node[aribox, right=of retire]  (dropr) {\figlabel{F15.retire_drop}};

  %% Loop enclosures: inner expansion sub-loop / outer parallel-run loop.
  \node[ariloop, fit=(pick)(prune)(expand)]                  (loopin)  {};
  \node[ariloop, fit=(batch)(react)(sterile)(append)(retire)] (loopout) {};
  \node[arinote, anchor=south west] at ([xshift=6mm]loopin.north)  (ttlin)  {\figlabel{F15.lane.inner}};
  \node[arinote, anchor=south west] at ([xshift=6mm]loopout.north) (ttlout) {\figlabel{F15.lane.outer}};
  \pic at ([xshift=-2.5mm]ttlin.west)  {ariglyph loop};
  \pic at ([xshift=-2.5mm]ttlout.west) {ariglyph loop};

  %% Step badges 1..7 (macro flow order; replaces edge labels).
  \node[aribadge] at (pick.north west)    {1};
  \node[aribadge] at (prune.north west)   {2};
  \node[aribadge] at (expand.north west)  {3};
  \node[aribadge] at (batch.north west)   {4};
  \node[aribadge] at (react.north west)   {5};
  \node[aribadge] at (sterile.north west) {6};
  \node[aribadge] at (retire.north west)  {7};

  %% Main top-down flow.
  \draw[ariarrow] (frontier) -- (pick);
  \draw[ariarrow] (pick)     -- (prune);
  \draw[ariarrow] (prune)    -- (expand);
  \draw[ariarrow] (expand)   -- (pending);
  \draw[ariarrow] (pending)  -- (batch);
  \draw[ariarrow] (batch)    -- (react);
  \draw[ariarrow] (react)    -- (sterile);
  \draw[ariarrow] (sterile)  -- (append);
  \draw[ariarrow] (append)   -- (retire);

  %% Gate-fired side effects: the one accent role in this figure.
  \draw[ariarrow accent] (prune)   -- (dropf);
  \draw[ariarrow accent] (sterile) -- (clamp);
  \draw[ariarrow accent] (retire)  -- node[arinote, above] {\figlabel{F15.yes_ab}} (dropr);

  %% Reflux: retire's "no" arm starts the next outer iteration.
  \draw[arifeedback] (retire.west) -- ([xshift=-1.6cm]retire.west)
                     |- (frontier.west);
  \node[arinote] at ([xshift=-1.6cm]pending.west) {\figlabel{F15.no_next}};
\end{tikzpicture}
}
  \caption{BFTS 单次外层迭代的状态机视图(自上而下的单一流程
           图)。步骤 1--3 为内层扩展子循环;步骤 4--7 为外层
           并行执行循环。黄色菱形为谓词门
           ($\prunepred,\sterilepred,\retirepred$);右侧红色
           虚线框为各谓词触发的副作用(前沿移除 / 分数夹紧 /
           retire)。图中未画出的辅助状态:每父节点的扩展计数
           器 $\expandcount(\cdot)$ 在步骤 3 自增,在步骤 7 的
           $\retirepred$ 中被使用;标签历史 $\Hist$(窗口
           $\histlen$)在步骤 5 由 \texttt{record\_run} 追加,
           在步骤 4 用于计算多样性奖励 $\divbonus{\cdot}$。
           可与 \cref{fig:bftsflow}(将相同循环呈现为带 LLM
           提示的细粒度控制流主线)对比。}
  \label{fig:bfts-state}
\end{figure}

\subsection{Run-loop 算法}\label{sec:exploration:run-loop}

\cref{alg:bfts-run-loop} 给出显式伪代码。内层 while 循环每次填
入一个空 worker 槽位(因此 worker 必在下一次提出扩展前开始运
行);外层块通过反复调用 \texttt{select\_next\_node}(每次一个节
点)组装批次直至达到 worker 上限,再并行执行该批次。执行后块
(a) 把已执行标签追加到 $\Hist$ 以使多样性奖励反映真实执行,
(b) 应用 $\retirepred$ 的规则 A / B。

\begin{algorithm}[t]
\small
\caption{BFTS run-loop(单次外层迭代)。}
\label{alg:bfts-run-loop}
\begin{algorithmic}[1]
\Require 状态 $\state=(\Frontier,\Pending,\expandcount,\Hist)$;
         配置 $(B_N, B_D, \maxexp)$ 与 worker 上限
         $\texttt{max\_parallel}=\min(\texttt{max\_parallel\_nodes},4)$
\Ensure  更新后状态 $\state'$
\State \Comment{--- 内层扩展子循环 ---}
\While{$\Frontier \neq \emptyset \land |\Pending| < \texttt{max\_parallel} \land |\Nodes| < B_N$}
  \State $f \gets \selectLLM(\Frontier)$ \Comment{\texttt{select\_best\_to\_expand}}
  \If{$\prunepred(f;\,|\Nodes|)$}
    \State $\Frontier \gets \Frontier \setminus \{f\}$;\;\;\textbf{continue}
  \EndIf
  \State $c \gets \texttt{expand}(f,\,\texttt{depth\_note},\,\texttt{budget\_note})$
  \State $\Pending \gets \Pending \cup \{c\}$;\;\;$\expandcount(f) \mathrel{+}{=} 1$
\EndWhile
\State \Comment{--- 外层并行执行批次 ---}
\State $B \gets \emptyset$
\While{$|B| < \min(\texttt{max\_parallel},\,|\Pending|)$}
  \State $n \gets \selectLLM(\Pending)$ \Comment{\texttt{select\_next\_node}:单个节点}
  \State $\Pending \gets \Pending \setminus \{n\}$;\;\;$B \gets B \cup \{n\}$
\EndWhile
\ForAll{$n \in B$ \textbf{并行}}
  \State $n' \gets \mathrm{ReAct}(n)$ \Comment{\texttt{AgentLoop.run},可能失败}
  \State $\Hist \gets (\Hist \mathbin{\Vert} \mathrm{label}(n'))[\,-\histlen{:}\,]$ \Comment{\texttt{record\_run}}
  \State $\Frontier \gets \Frontier \cup \{n'\}$
  \ForAll{$p \in \Frontier : p = \parentof{n'}$} \Comment{规则 A}
    \If{$\nodescore{n'} > \nodescore{p}$}
      \State $\Frontier \gets \Frontier \setminus \{p\}$
    \EndIf
  \EndFor
  \ForAll{$p \in \Frontier$} \Comment{规则 B}
    \If{$\expandcount(p) \geq \maxexp$}
      \State $\Frontier \gets \Frontier \setminus \{p\}$
    \EndIf
  \EndFor
\EndFor
\State \Return $\state' = (\Frontier,\Pending,\expandcount,\Hist)$
\end{algorithmic}
\end{algorithm}

\subsection{节点选择(LLM 主导)}\label{sec:exploration:node-selection}

\begin{figure}[t]
  \centering
  \resizebox{\columnwidth}{!}{% F03: signals consumed by the LLM-driven node selector.
% Replaces the earlier closed-form utility figure: ARI does not pin
% a scalar utility; instead a structured candidate description is
% handed to the LLM and an index is parsed back. See F11 for the full
% control flow; this figure focuses on the *information* that flows in.
% style.tikzstyles is loaded by shared/preamble.tex / standalone wrapper.
\begin{tikzpicture}[node distance=0.5cm and 1.6cm]

  % per-node signals (left column), all stations for paper-style typography
  \node[station]                                       (real)   {\figlabel{F3.has_real}};
  \node[station, below=of real]                        (metrics){\figlabel{F3.metrics}};
  \node[station, below=of metrics]                     (depth)  {\figlabel{F3.depth_signal}};
  \node[station, below=of depth]                       (label_st){\figlabel{F3.retry}};
  \node[station, below=of label_st]                    (summ)   {\figlabel{F3.summary}};
  \node[station, fill=ari-pink!22, below=of summ]      (div)    {\figlabel{F3.div}};

  % LLM selector (centre) — wider station emph
  \node[station emph big, right=2.2cm of metrics, yshift=-1.5cm]
                                                       (sel)    {\figlabel{F3.selector}};
  % the consumed prompt file (above the selector)
  \node[station, fill=ari-yellow!30, draw=ari-orange,
        above=0.6cm of sel]                            (prom)   {\figlabel{F3.prompt}};

  % output (right) — single index back to caller
  \node[station emph big, right=2.2cm of sel]          (idx)    {\figlabel{F3.index}};

  % deterministic fallback when LLM index is unparseable
  \node[station, dashed, draw=ari-vermilion, below=1.0cm of sel]
                                                       (fb)     {\figlabel{F3.fallback}};

  % --- numbered step badges --------------------------------------
  \node[numbered step, above left=0.05cm and -0.2cm of real.north west] {1};
  \node[numbered step, above left=0.05cm and -0.2cm of sel.north west]  {2};
  \node[numbered step, above left=0.05cm and -0.2cm of idx.north west]  {3};

  % --- edges -----------------------------------------------------
  \foreach \n in {real,metrics,depth,label_st,summ}
    \draw[arrow] (\n.east) -- (sel.west);
  \draw[arrow dashed] (div.east) to[bend right=8] (sel.west);
  \draw[arrow dashed] (prom.south) -- (sel.north);
  \draw[big arrow] (sel.east) -- (idx.west)
       node[edge label, font=\sffamily\footnotesize]{\figlabel{F3.parse}};
  \draw[big arrow, draw=ari-vermilion] (sel.south) -- (fb.north)
       node[edge label, font=\sffamily\footnotesize]{\figlabel{F3.unparseable}};
  \draw[big arrow, draw=ari-vermilion] (fb.east) to[bend right=22] (idx.south);

  % --- background grouping for input column ----------------------
  \begin{scope}[on background layer]
    \node[stage container, fit=(real)(metrics)(depth)(label_st)(summ)(div)] (col) {};
    \node[stage title, above=4pt of col.north west, anchor=south west] {\figlabel{F3.band.signals}};
  \end{scope}
\end{tikzpicture}
}
  \caption{交给 LLM 用以挑选下一节点的选择信号。输入被拼接成结
           构化候选描述;软性的 \texttt{diversity\_bonus} 仅作平
           手解决,不作主信号。}
  \label{fig:nodescore}
\end{figure}

ARI 有意\emph{不}将节点排序归约为闭式标量效用。相反,
\texttt{select\_next\_node}(「下一智能体步该运行哪个 pending 节
点」)与 \texttt{select\_best\_to\_expand}(「哪个完成节点最值得
扩展」)都把候选描述列表交给 LLM,并解析返回的单个 0 起始索引
—— 每次调用恰好选出一个节点。每个节点 $n$ 的候选描述形如:
\begin{quote}\small\texttt{[i] id=..., depth=..., label=..., has\_real\_data=..., metrics=\{...\}, summary=...}\end{quote}
运行选择侧(\code{select\_next\_node})会对当前获得奖励的节点
在行尾追加 \code{diversity\_bonus=+0.05};\code{select\_best\_to\_expand}
不追加。消费它的提示分别于 \cref{prompt:bfts_select}(选择)与
\cref{prompt:bfts_expand_select}(扩展目标)逐字载录。提示中列
出的选择标准:\texttt{has\_real\_data=True} 且
metrics 强者高价值;metrics 整体多目标加权;深而成果好的节点值
得继续;\texttt{diversity\_bonus} 仅当软性平手解决项处理。
\texttt{label} 字段与展示给 \texttt{select\_best\_to\_expand} 的
信息一致,并且不存在误导性的"偏好低 retry"标准。

\subsubsection*{多样性奖励}\label{sec:exploration:diversity}

多样性奖励 \texttt{BFTS.diversity\_bonus} 以
\texttt{\_recent\_label\_history}(由 \texttt{record\_run} 填充
的滑动窗)跟踪近期标签。对于其标签在该窗中相对最常见标签为
\emph{欠表达}的节点,函数返回:
\begin{equation}\label{eq:diversity}
  \mathrm{div}(n) =
    \begin{cases}
      +0.05 & \mathrm{cnt}(n) \le \tfrac{1}{2}\,\mathrm{cnt}_{\max},\\
      0     & \text{否则,}
    \end{cases}
\end{equation}
其中 $\mathrm{cnt}(n)$ 为窗内 $\mathrm{label}(n)$ 的出现次数,
$\mathrm{cnt}_{\max}=\max_\ell\mathrm{cnt}(\ell)$。
$0.05$ 这个常数有意小:scientific score 仍主导排序,奖励仅打破
近似平手。

\subsubsection*{LLM 回退}\label{sec:exploration:fallback}

若 LLM 返回无法解析的索引,\code{select\_next\_node} 通过
\code{\_select\_fallback} 辅助函数(按 \code{\_fallback\_score}
式对候选排序)回退至确定性排序。默认评分式为
\begin{equation}\label{eq:fallback}
  \fallbackscore{n} \;=\; \nodescore{n} + \mathrm{div}(n),
\end{equation}
其中 $\nodescore{n}\equiv\texttt{\_scientific\_score}(n)$ 是节点
的标量奖励(\cref{sec:notation})。\texttt{\_select\_fallback}
在存在 \texttt{has\_real\_data=True} 节点时优先选取它们,然后返
回 $\fallbackscore{\cdot}$ 最大的候选。即使 LLM 调用退化,循环
也保证向前推进。

\subsubsection*{评估层的配置}\label{sec:exploration:config-layers}

\cref{sec:exploration:formalization,sec:exploration:node-selection,sec:exploration:fallback}
中出现的 4 个评估面可独立配置(默认值复刻上式;未修改的配置等同
于 no-op)。

\begin{description}
  \item[\textbf{合成公式.}]
        由 \code{evaluator.composite} 选择。
        $\{a_k\}\!\to\!\texttt{\_scientific\_score}$ 的合成方式可选:
        加权调和平均(默认)、加权算术平均、瓶颈式
        \code{weighted\_min}、加权几何平均。调和与几何平均使用
        $\epsilon=0.01$ 对零值轴下界,保证分母有限。
  \item[\textbf{前沿评分策略.}]
        由 \code{bfts.frontier\_score} 选择。
        \cref{eq:fallback} 对应 \code{scientific\_plus\_diversity};
        其余可选 \code{scientific\_only}(去除 $\mathrm{div}$);
        \code{depth\_penalized} 在 $\mathrm{div}(n)$ 之上叠加
        $-\lambda\,\depth{n}$,其中
        $\lambda=$\code{depth\_penalty\_lambda};
        \code{ucb\_like} 在 $\mathrm{div}(n)$ 之上叠加
        $c_{\mathrm{ucb}}\sqrt{\log N / (\expandcount(n)+1)}$,其中
        $c_{\mathrm{ucb}}=$\code{ucb\_c}、$N = \sum_{n'}\expandcount(n')+|\Frontier|$。
  \item[\textbf{评估轴集合.}]
        由 \code{evaluator.axis\_mode} 选择。
        \texttt{dynamic}(默认)= 通用 5 轴 floor + 有效 rubric +
        \texttt{idea.json} 的 plan-keyword 轴;\texttt{legacy} 固定
        5 轴;\texttt{custom} 直接使用
        $\{(\mathrm{name},\mathrm{description},w)\}$ 列表。
  \item[\textbf{选择 prompt.}]
        由 \code{bfts.select\_prompt} 与
        \code{bfts.expand\_select\_prompt} 设置。
        通过 \code{FilesystemPromptLoader} 键替换
        \cref{alg:bfts-run-loop} 中两个 $\selectLLM$ 所用模板。
        选择模板须接受占位符 \code{experiment\_goal}、
        \code{memory\_context} 与 \code{candidates};扩展目标模板接受
        \code{experiment\_goal} 与 \code{candidates}。
\end{description}

\subsection{扩展(每次调用 1 个子节点)}\label{sec:exploration:expansion}

\texttt{expand} 函数每次调用\emph{至多}生成一个新子节点。同一
父节点希望多个子节点的调用方反复调用 \texttt{expand},通过
\texttt{existing\_children} 参数把已生成的子节点串入,使 LLM 避
免提议重复方向。

各新子的标签由 LLM 判断而非硬编码模板决定;扩展提示(逐字载录
于 \cref{prompt:bfts_expand})接收:

\begin{itemize}
  \item 父节点状态(\texttt{succeeded} /
        \code{failed/no-real-data})与
        \code{\_scientific\_score};
  \item 父节点结构化的 \texttt{node\_report}
        (\code{delta\_vs\_parent} / \texttt{concerns} /
        \texttt{hints},由 node-report 构建器生成)若存在,以便
        规划者瞄准具体弱点;
  \item 同深度的兄弟得分,以及自根至父的祖先得分;
  \item 整树多样性指标(目前已见的唯一标签、深度分布);以及
  \item 此父节点已生成的子节点标签,避免重复提议。
\end{itemize}

子节点以 \StateOpen{} 创建,交给智能体循环执行
(\cref{sec:exploration:react}),所有轴填齐后变为
\StateDone{};智能体循环抛异常时变为 \StateFailed{}。

\subsection{剪枝与终止}\label{sec:exploration:termination}

\texttt{should\_prune} 实现剪枝谓词 $\prunepred$(\cref{eq:prune}):
总节点上限、深度上限与 sterile 标志。深度检查激活了
\texttt{max\_depth} 配置。前沿退役谓词 $\retirepred$
(\cref{eq:retire})在每个节点完成后施加:规则 A 在子的奖励
$\nodescore{\cdot}$ 超过父 $f$ 时退役 $f$,规则 B 在 $f$ 已被扩展
$\maxexp$ 次时退役。$\retirepred$ 不删除节点,其历史仍保留在
$\Tree$ 中,只是收缩前沿池,使搜索分散而不是无限挖掘同一父节点。

下列任一条件满足即停止:

\begin{itemize}
  \item 全局 \texttt{max\_total\_nodes} 达到;
  \item 前沿中所有节点均触发 $\Pi$ 且 pending 为空;
  \item 调用侧成本追踪器强制的步数 / 成本预算耗尽;
  \item \term{停滞检测器} \texttt{detect\_stagnation} 观察到
        \texttt{\_scientific\_score} 的运行最大值在滑动窗内不再
        改善;
  \item 谱系判断模块的 \texttt{LineageDecision} 显式终止该分支
        (参见 \cref{prompt:lineage_decision})。
\end{itemize}

这些当中实际最常发生哪一种,是一个经验问题,留待冻结检查点就位
后由 \cref{chapter:eval} 回答;本报告不主张其分布。

\subsection{谱系与可复现性}\label{sec:exploration:lineage}

谱系即最佳节点祖先的链条,作为 \texttt{tree.json} 中的
\texttt{lineage} 键写入检查点。\texttt{LineageState} 数据类捕
获 BFTS 用于决定是否继续的滚动统计;跨检查点的遍历器
\texttt{walk\_ancestor\_ckpts} 提供跨运行的父指针。构想池经
\texttt{get\_idea\_pool\_for\_ckpt} 继承,排序后的根构想选择记
录在 \texttt{RootChoice} 数据类中以便事后审计(提示词:
\cref{prompt:root_idea_selector})。

可复现性依赖三项不变量。其一,每次随机化操作以节点标识符播种。
其二,每次工具调用将参数与输出写入所属节点的工作目录,绝不写
入父节点 —— 由 MCP 层的 \texttt{cow\_node\_id} 校验强制。其
三,成本账本在 \texttt{CostTracker.init} 处为每节点附加美元金
额,因此给定相同提示,预算检查保持确定。

\paragraph{默认离线。}
可复现性也规定搜索循环可以读取什么。默认下每个节点的智能体循环
\emph{离线}运行 —— Web 访问技能被相位限制到论文撰写与复现阶段,
因此节点无法在搜索过程中查询随时变化的实时检索结果,每节点轨迹
也就不依赖于其运行当天的 Web 状态。文献查阅仍会发生,但更早且在
循环之外 —— 在搜索开始前、构想阶段执行的一次有界综述。确实需要
实时信息的运行可以选择性开启(\code{bfts.allow\_web}),在探索期
间向节点智能体开放 Web 访问;此时运行会向其检查点写入一个不可复
现轨迹标记,使日后的复现不会被默默当作精确。无论哪种方式,确定
性契约都保持诚实:默认运行按字节可复现,而以可复现性换取实时信
息的运行会把这一权衡记录在自身来历中。

\subsection{面向计算型证据的谱系串联}\label{sec:exploration:lineage-chaining}

并非运行打算支撑的每条主张都来自新的测量。有些证据是从已经存在
的测量\emph{计算}而来——对早先探针数据的参数拟合、对该拟合的留
出验证(held-out validation)、在多个构型间基于模型的选型。事实
证明,这类主张对相互独立的树节点是结构性不可达的:在一次真实运
行中,以拟合或验证为目标、却从不携带测量的父节点扩展出的子节点,
退化为重新运行单个探针——因为「用你已有的数据来计算」从来不是一
个可见的选项:子节点手里没有数据,也没有任何信号告诉选择器哪个
节点有。

弥合这一缺口的机制只使用树本来就有的父节点$\to$子节点工作目录
继承(子节点在父节点工作树的副本内执行;\cref{sec:paper:evidence}
依赖同一事实),在扩展的两侧各加一个提示。父侧:扩展目标选择的
提示(\cref{sec:exploration:node-selection})被扩充一个覆盖度提
示,点名其工作树持有最多\term{契约证据名}的节点——即本次运行已
声明的指标契约(最终论文必须举证的可证伪主张集合;
\cref{sec:paper:contract} 给出完整定义)所要求的测量名——使拟合
或验证型的子节点优先从已持有其输入的节点扩展。子侧:当子节点继
承的目录中含有谱系测量文件时,它在起始即被告知哪些文件存在、其
中含有哪些确切的契约名,并被指示从这些文件计算派生证据,而非重
新运行底层实验。

跨越节点边界的只有名称:文件名与测量名,从不是数值,也从不是兄
弟分支的结论。这些提示与契约本身同属运行级协调元数据,因此
\cref{sec:exploration:memory} 的分支隔离原样保持——有缺陷分支的
数字仍然无法泄漏进兄弟分支的推理——同时计算型证据链(先拟合、再
验证、后选型)得以沿单一谱系执行,而不是塌缩回反复的探针运行。

\subsection{ReAct 驱动器}\label{sec:exploration:react}

\begin{figure}[t]
  \centering
  \resizebox{\columnwidth}{!}{% F04: ReAct loop — Think / Act / Observe / Reflect with numbered stages,
%       paper-style typography.
% style.tikzstyles is loaded by shared/preamble.tex / standalone wrapper.
\begin{tikzpicture}[node distance=2.0cm and 2.5cm]

  \node[state init]                            (think) {\figlabel{react.think}};
  \node[state, right=of think]                 (act)   {\figlabel{react.act}};
  \node[state, below=of act]                   (obs)   {\figlabel{react.observe}};
  \node[state final, below=of think]           (refl)  {\figlabel{react.reflect}};

  % numbered step badges, perched on each state's upper-left corner
  \node[numbered step, above left=0.05cm and -0.15cm of think.north west] {1};
  \node[numbered step, above left=0.05cm and -0.15cm of act.north west]   {2};
  \node[numbered step, above left=0.05cm and -0.15cm of obs.north west]   {3};
  \node[numbered step, above left=0.05cm and -0.15cm of refl.north west]  {4};

  % external annotations (yellow = memory, orange = phase-filtered tools)
  \node[station, fill=ari-orange!22, draw=ari-vermilion,
        above=0.6cm of act]                    (tools) {\figlabel{F4.toolset}};
  \node[station, fill=ari-yellow!30, draw=ari-orange,
        left=0.6cm of think]                   (mem)   {\figlabel{F4.memcontext}};

  % --- big arrows (paper-style) ---
  \draw[big arrow] (think) -- (act)   node[edge label, font=\sffamily\footnotesize]{\figlabel{F4.tool_call}};
  \draw[big arrow] (act)   -- (obs)   node[edge label, font=\sffamily\footnotesize]{\figlabel{F4.execute}};
  \draw[big arrow] (obs)   -- (refl)  node[edge label, font=\sffamily\footnotesize]{\figlabel{F4.parse}};
  \draw[big arrow] (refl)  -- (think) node[edge label, font=\sffamily\footnotesize]{\figlabel{F4.update}};
  \draw[arrow dashed] (obs) to[bend left=18] (think);

  \draw[arrow dashed] (tools.south) -- (act.north);
  \draw[arrow dashed] (mem.east)    -- (think.west);

  % budget annotation
  \node[stage title, below=0.3cm of refl] {\figlabel{F4.maxsteps}};
\end{tikzpicture}
}
  \caption{ReAct 循环。实线为主循环;虚线反向边表示当观测已闭
           合开放问题时的提前退出。}
  \label{fig:react}
\end{figure}

每个节点的扩展在 \texttt{AgentLoop} 类中运行 ReAct 风格的智能
体循环。最大步数由 \texttt{MAX\_REACT\_STEPS = 80} 控制,可按
实例通过 \texttt{max\_react\_steps} 关键字覆盖。另一保护
\texttt{MIN\_TOOL\_CALLS = 2} 防止平凡过短的轨迹。智能体的系统
提示词逐字载录于 \cref{prompt:agent_system}。

智能体可用工具集合按阶段由工具管理器过滤;例如探索阶段屏蔽论
文撰写工具,使 BFTS 扩展不会捷径地变成草稿。一小部分 MCP 工具
为 \texttt{ari-core} 自身预留(memory CoW 操作)并对 LLM 隐
藏;这保证模型无法设置任意 node id 来绕过记忆技能的写时复制校
验。

\subsection{研究记忆}\label{sec:exploration:memory}

节点并非孤立运行:每个节点继承其祖先所确立的内容,反过来又为其
后代与论文撰写器留下记录。ARI 把这一记录保存在每次运行的\term{研
究记忆}中,它遵循自由形式日志无法满足的两项要求。其一,\emph{分
支隔离}:节点只能读取自身祖先的记忆,绝不读取兄弟或无关检查点的
记忆,因此两条并行线索无法相互污染 —— 与
\cref{sec:exploration:lineage} 中定义谱系的祖先作用域谓词相同。
其二,\emph{可基底性}:由于记忆最终支撑论文,条目必须能指向结果
背后的证据,而非仅复述结果,从而一个主张可回溯到支撑它的产物。

\subsubsection*{作为证据索引的记忆}

节点产出内容的唯一真值来源是其结构化自评报告 —— 节点完成时写出
的 \texttt{node\_report.json}(\cref{sec:paper:evidence} 详述其字
段)。记忆条目不复制这些字段。它携带一段简短可检索文本、一个
\term{种类}、一个指回原始节点报告的指针,以及对结果所依赖的特定
产物 —— 日志、源文件、输出文件 —— 的内容哈希引用。因此记忆条目是
\emph{对证据的索引},而非证据本身:它所引用的产物可随时重新哈
希,以确认其仍存在且仍匹配,这一完整性检查把每个引用归类为已验
证、缺失或不匹配。\term{种类}取自固定词汇 —— 观测、实验结果、失
败案例、过程、反思、产物摘要、论文主张、可复现性事件 —— 使读者能
精确索取所需记录:比如某条分支上的失败,或仅由产物支撑的结果。

\subsubsection*{节点终端的整合}

节点完成时,一个确定性的\term{整合}步骤读取其节点报告并发出带类
型条目:为产生测量的节点发出实验结果条目,附带对其所测产物的来
历引用;为失败的节点发出失败案例条目。该步骤不调用语言模型,因
此同一节点报告总是整合为相同条目 —— 整合是已记录状态的纯函数,
重跑检查点会再现它所产生的记忆。它是循环钩子而非智能体动作:搜
索智能体从不会被要求“保存到记忆”,实际上智能体也不会自行去取用
回忆工具。循环所依赖的记忆读取由循环自身确定性地完成 —— 正是这
一点让记忆能承载来历,而不必把非确定性的检索放在关键路径上
(\cref{sec:discussion:p2})。

\begin{figure}[t]
  \centering
  \resizebox{0.92\columnwidth}{!}{% F18: verifiable research memory as a TYPED INDEX OVER EVIDENCE (not a flow).
% The store holds typed records that point, by content hash, at the artefacts
% that ground them; reproduction outcomes are appended and folded to a latest
% status; an admissibility predicate usable(mu) gates which records may become
% paper claims, so the SAME store yields two projections: a working context for
% exploration and a verified context for the paper.
% style.tikzstyles is loaded by shared/preamble.tex / the standalone wrapper.
\begin{tikzpicture}[node distance=1.5cm and 2.2cm]

  % --- source of truth feeding the store -----------------------------
  \node[artifact wide]                          (report) {\figlabel{F18.report}};
  \node[block emph big, below=1.3cm of report]  (store)  {\figlabel{F18.store}};

  % --- what the records index, and what they accrete -----------------
  \node[block native big, left=of store]        (artifacts) {\figlabel{F18.artifacts}};
  \node[block native big, right=of store]       (repro)     {\figlabel{F18.repro}};
  % record vocabulary: a legend in the empty space under the artefacts
  \node[block aux, below=1.3cm of artifacts]    (kinds)     {\figlabel{F18.kinds}};

  % --- two projections of the SAME store -----------------------------
  \node[block native big, below=2.4cm of store]              (workctx) {\figlabel{F18.workctx}};
  \node[guard, below right=1.4cm and 1.7cm of store]         (usable)  {\figlabel{F18.usable}};
  \node[block emph big, below=1.1cm of usable]               (vctx)    {\figlabel{F18.vctx}};

  % --- edges ----------------------------------------------------------
  \draw[arrow]        (report) -- node[edge label] {\figlabel{F18.consolidate}} (store);
  \draw[arrow dashed] (store)  -- node[edge label] {\figlabel{F18.refs}}        (artifacts);
  \draw[arrow]        (repro)  -- node[edge label] {\figlabel{F18.fold}}        (store);
  \draw[arrow dashed] (store)  -- (kinds);
  \draw[arrow]        (store)  -- node[edge label] {\figlabel{F18.forexpl}} (workctx);
  \draw[arrow]        (store)  -- (usable);
  \draw[arrow]        (usable) -- node[edge label] {\figlabel{F18.forpaper}} (vctx);

\end{tikzpicture}
}
  \caption{可验证研究记忆。对每个节点真值来源
           \texttt{node\_report.json} 的确定性整合,填充一个按内容
           哈希\emph{索引}产物(日志、输出、源码)的\emph{带类型}
           存储。复现事件被追加并折叠为最新状态,准入谓词
           $\mathrm{usable}(\mu)$(\cref{eq:usable})控制哪些记录
           进入已验证上下文。同一存储供养两个读者:在下一节点起始
           注入的\emph{工作上下文}(探索的继承路径),以及为论文定量
           主张提供依据的\emph{已验证上下文}
           (\cref{sec:paper:evidence})。无据的反思可引导探索,但绝
           不引导论文正文。}
  \label{fig:research-memory}
\end{figure}

\subsubsection*{两个读者:工作上下文与已验证上下文}

该存储供养两个消费者(\cref{fig:research-memory}),二者之别正是
设计的核心。在每个节点\emph{起始},循环注入\term{工作上下文}:实
验的固定核心 —— 目标、目标度量、硬件 —— 以及其祖先所记录的结论。
于是子节点不是从聚合文本转储开始,而是从谱系已确立的内容开始,确
定性地且限定于自身分支。在\emph{论文}阶段,流水线从最佳节点的谱
系构建\term{已验证上下文}:它把每个结果的可复现性事件折叠为最新
状态,并对条目排序,使可复现的结果高于仅由产物支撑者,后者又高
于无据的反思。条目 $\mu$ 仅当其由产物支撑且未被复现尝试反驳时,
才\emph{可作为}定量论文主张:
\begin{equation}\label{eq:usable}
  \mathrm{usable}(\mu) \;=\;
    \mathrm{grounded}(\mu)\;\wedge\;
    \bigl(\mathrm{repro}(\mu)\neq\mathrm{failed}\bigr).
\end{equation}
无据的反思仍可通过工作上下文引导搜索,但 \cref{eq:usable} 把它挡
在论文正文之外:撰写器的定量主张只依赖于可作为主张、优先可复现的
条目。这是证据组装器(\cref{sec:paper:evidence})在记忆一侧的对
应物;组装器决定撰写器可读取哪些产物\emph{字节},而已验证上下文
决定哪些\emph{结果}可成为主张。

\subsubsection*{作为追加事件的可复现性}

由于过去条目按字节稳定 —— 节点只在自身标识符下写入,既有条目从
不就地修改 —— 结果的可复现性状态在日后重跑时不会被就地改写。重
跑改为\emph{追加}一个可复现性事件条目,任何读者都会把某目标的事
件折叠为其最新状态。复现失败的结果因此被从 \cref{eq:usable} 的主
张集合中降级 —— 并可诚实地作为局限报告 —— 而不必改写当初把它记
录为有望的历史。
